% Sean P. Florez | resume
% Layout and typography follow Sean's supplied LaTeX template.
% Compile: pdflatex main.tex
\documentclass[letterpaper,10pt]{article}
\usepackage[left=0.65in,right=0.65in,top=0.50in,bottom=0.55in]{geometry}
\usepackage{enumitem}
\usepackage[colorlinks=true,urlcolor=blue]{hyperref}
\usepackage{titlesec}
\usepackage{parskip}
\usepackage{amsmath}
\usepackage{needspace}

% Searchable text; the default Computer Modern typography is unchanged.
\input{glyphtounicode}
\pdfgentounicode=1
\hypersetup{
    pdftitle={Sean P. Florez | Resume},
    pdfauthor={Sean P. Florez}
}

% Section spacing and rule, as supplied.
\setlength{\parindent}{0pt}
\setlength{\parskip}{3pt}
\titlespacing\section{0pt}{9pt plus 2pt minus 1pt}{5pt plus 1pt minus 1pt}
\titleformat{\section}{\large\bfseries}{\thesection}{1em}{}[\titlerule]

% List style, as supplied.
\setlist[itemize]{leftmargin=*,noitemsep,topsep=0pt}
\setlength{\emergencystretch}{1em}
\widowpenalty=10000
\clubpenalty=10000
\raggedbottom

% Page-break guards: keep a heading with the lines that follow it.
\newcommand{\keep}{\needspace{4\baselineskip}}
\newcommand{\keepsection}{\needspace{6\baselineskip}}

\begin{document}
\pagestyle{empty}

% ===== Header =====
\begin{center}
    {\LARGE \textbf{Sean P.\ Florez}}\\
    \href{mailto:sean.florez@colorado.edu}{sean.florez@colorado.edu} \,\textbar\, (786) 202--5558 \,\textbar\, Boulder, CO\\
    \href{https://seanflorez.com}{seanflorez.com} \,\textbar\,
    \href{https://linkedin.com/in/sean-florez}{linkedin.com/in/sean-florez} \,\textbar\,
    \href{https://github.com/fl-sean03}{github.com/fl-sean03}
\end{center}

\vspace{-2pt}
\begin{center}
\textit{Security Clearance: Top Secret (active); SCI eligible.}
\end{center}

\vspace{-6pt}
\begin{center}
\textit{Materials R\&D \;\textbar\; Research Automation \;\textbar\; Manufacturing Systems}
\end{center}

% ===== Synopsis =====
\keepsection\section*{Synopsis}
Materials science Ph.D.\ researcher building \textbf{materials modeling tools}, \textbf{research automation}, and \textbf{manufacturing qualification systems}. Experience running large computational campaigns, supporting a new energetic-materials production line, and evaluating technologies for DoD and DOE.

% ===== Materials research =====
\keepsection\section*{Materials Research \& Federal Laboratories}
\textbf{Graduate Research Assistant} \hfill \textit{University of Colorado Boulder, Oct 2024--Present}\\
\textit{Heinz Interfaces Laboratory, Prof.\ Hendrik Heinz}
\begin{itemize}
    \item Developed metal-oxide force-field parameterization for an \textbf{industry-funded force-field program}, using experimental lattice constants, elastic constants, surface energies, and vibrational spectra.
    \item Validated a Cu\(_2\)O parameter family within \textbf{5\%} of both experimental elastic references; verified \textbf{52 configurations} for LAMMPS, msi2lmp, CHARMM, and NAMD.
    \item Built structural-neighbor retrieval over \textbf{3,009 structures and 6.7 million atoms} to assign starting atom types and parameters, with confidence thresholds and explicit coverage gaps.
    \item Executed \textbf{780 MXene shear trajectories} across six termination and geometry conditions with \textbf{zero simulation failures}; used a frozen protocol and independent audit of raw trajectories.
    \item Ran \textbf{60 Pt-facet production simulations} and nanocrystal ensembles to investigate hydrogen-release configurations. Wrote the technical narrative for an ALCF award of approximately \textbf{26,150 node-hours}.
\end{itemize}

\keep\textbf{High-Performance Computing Intern} \hfill \textit{Air Force Research Laboratory, May--Aug 2025}\\
\textit{Materials \& Manufacturing Directorate, Wright-Patterson AFB}
\begin{itemize}
    \item Built LAMMPS shear workflows for hydrated Ti\(_3\)C\(_2\) MXene bilayers and ran parameter sweeps on \textbf{DoD HPC}.
    \item Quantified hydration-dependent interlayer shear at 103 MPa dry, 8 MPa at quarter-monolayer coverage, and 39 MPa at full coverage.
\end{itemize}

\keep\textbf{Materials Modeling Intern} \hfill \textit{Air Force Research Laboratory, May--Dec 2023}
\begin{itemize}
    \item Automated high-throughput DFT screening of low-work-function emitters using Quantum ESPRESSO, ASE, and pymatgen.
\end{itemize}

\keep\textbf{Undergraduate Research Assistant} \hfill \textit{Hennig Group, University of Florida, Oct 2021--May 2024}
\begin{itemize}
    \item Benchmarked VASPsol against experimental solvation energies; optimized cavity and screening parameters with grid search, Nelder--Mead, and surrogate modeling.
\end{itemize}

% ===== Research automation =====
\keepsection\section*{Research Automation \& Compute Infrastructure}
\textbf{OpenSDL, computational and autonomous laboratory framework} \hfill \textit{2026--Present}\\
\href{https://github.com/fl-sean03/OpenSDL}{github.com/fl-sean03/OpenSDL} \;\textbar\; Apache-2.0; documentation, conformance tests, and CI
\begin{itemize}
    \item Built typed operations with resource requirements, risk classes, timeouts, and retries; supported interchangeable robot, simulator, and human executors.
    \item Implemented campaign scoring and next-round selection. A reference three-dye formulation search matched an unseen target recipe within \textbf{0.5 percentage points}, evaluating \textbf{96 candidates per round}.
    \item Implemented resource leases, restart recovery, policy checks, and portable run records; developed adapters for a simulated lab, local compute, and human tasks, with Slurm and hardware integration in progress.
\end{itemize}

\keep\textbf{agent-fleet, persistent research automation} \hfill \textit{2026--Present}\\
\href{https://github.com/fl-sean03/agent-fleet}{github.com/fl-sean03/agent-fleet}
\begin{itemize}
    \item Built a research harness based on a \textbf{15-agent fleet} operating for months, with isolated workspaces, persistent conversations, and inter-agent messaging.
    \item Ran multi-day computational campaigns with recovery after reboots, credential expiry, and preemption.
\end{itemize}

\keep\textbf{HPC \& Cloud Throughput Infrastructure} \hfill \textit{2025--Present}
\begin{itemize}
    \item Developed \textbf{subjob} to multiplex short tasks within one HPC allocation, with shared worker pools and heartbeat recovery.
    \item Built \textbf{CloudComputeManager} for GPU spot-instance provisioning, environment setup, pipeline execution, preemption recovery, and per-run cost tracking; used it for MXene shear sweeps.
\end{itemize}

% ===== Manufacturing =====
\keepsection\section*{Manufacturing \& Production Systems}
\textbf{Engineering Contractor, Rocky Mountain Scientific Laboratory} \hfill \textit{Jun--Aug 2026}\\
\textit{Quality system for a client affiliate's new energetic-materials production line}\\
\textit{Automation \& Robotics embed, Jun--Aug 2026}
\begin{itemize}
    \item Developed lot-release records covering inspections, approval authority, hold criteria, and nonconformance disposition, with item-level traceability to the governing \textbf{military specification}.
    \item Built a controlled, single-source document set with \textbf{10 client-format instructions}, \textbf{20 test-specific quality-control packets}, manufacturing procedures, and inspections.
    \item Automated checks for process-route isolation, link validity, terminology, and consistency between released documents and their source.
    \item Built a program tracker that generates the Kanban board, decision log, requirements traceability matrix, and design-review record.
    \item Supported quality, packaging, labeling, on-site commissioning, and hazard and failure-mode analysis.
\end{itemize}

% ===== Government =====
\keepsection\section*{National Security \& Technology Transition}
\textbf{S\&T Scouting Intern} \hfill \textit{The Joint Staff, J7 / Future Technology Office, Aug--Oct 2025}
\begin{itemize}
    \item Produced \textbf{eight technical briefs} on propulsion and energy concepts, assessing readiness, risk, and experimentation paths for TRL 3--5 technologies across Service laboratories and DoD R\&D organizations.
\end{itemize}

\keep\textbf{Technology Commercialization Intern} \hfill \textit{Idaho National Laboratory (DOE), May--Aug 2024}
\begin{itemize}
    \item Conducted \textbf{20 Energy I-Corps interviews}, developed a commercialization plan for a laboratory technology, and supported a DOE proposal awarded \textbf{\$2M}.
\end{itemize}

% ===== Education =====
\keepsection\section*{Education}
\textbf{University of Colorado Boulder} \hfill \textit{Expected May 2028}\\
Ph.D., Materials Science \& Engineering. Advisor: Prof.\ Hendrik Heinz.

\textbf{University of Florida} \hfill \textit{May 2024}\\
B.S., Materials Science \& Engineering. Minor in Chemistry.

% ===== Skills =====
\keepsection\section*{Core Skills}
\textbf{Materials \& Modeling:} LAMMPS; NAMD; VASP/VASPsol; Quantum ESPRESSO; ASE; pymatgen; MDAnalysis; force-field parameterization; elastic and surface-energy validation.\\
\textbf{Research Automation:} Campaign orchestration; resource scheduling; adapters for a simulated lab, local compute, and human tasks; run provenance; restart recovery; experiment-selection workflows; policy checks.\\
\textbf{Software \& HPC:} Python; Bash; Linux; Git; Docker; SQLite; pytest; CI; Slurm; PBS Pro; checkpoint/restart; GPU and CPU-MPI workflows.\\
\textbf{ML \& Data:} scikit-learn; TensorFlow; RDKit; SOAP descriptors; regression and classification; surrogate modeling; uncertainty analysis.\\
\textbf{Manufacturing \& Quality:} Lot-release records; military-specification traceability; controlled procedures; nonconformance disposition; hazard and failure-mode analysis; bills of materials; vendor qualification.

% ===== Research outputs =====
\keepsection\section*{Selected Research Outputs}
\textbf{Co-author:} \textit{INTERFACE Force Field for Alumina with Validated Bulk Phases and a pH-Resolved Surface Model Database for Electrolyte and Organic Interfaces.} \href{https://arxiv.org/abs/2601.12570}{arXiv:2601.12570}, 2026 (preprint).

\textbf{Lead-author manuscripts in preparation:} \textit{Modeling of MXene Shear Properties via Atomistic Molecular Dynamics Simulations}, with AFRL Materials \& Manufacturing Directorate co-authors;\\ \textit{Understanding Hydrogen Release Mechanisms of N-Ethylcarbazole on Pt Nano-Crystal Catalysts}.

\end{document}
